%% notes-tex/024-perturb-seq-pilot/sections/2-routes.tex
%% SOURCE: hand-authored synthesis of notes-tex/024-perturb-seq-costing/
%%         sections/3-method-explainers.tex (Secs. 3.1, 3.2, 3.7) and
%%         sections/5-application.tex (Sec. 5.5, the core meeting). No new
%%         numbers; each is cross-referenced to the section that derives it.

\section{The Three Bench Routes\secstatus{todo}}
\label{sec:routes}

All three routes end in the same object, a table of cells by genes with a
perturbation label on each cell. They differ in how a cell keeps its identity
while its RNA is being converted to a sequencing library, and that difference
decides who does the work. Two of them are the familiar pair, plate barcoding
and droplet isolation. The third composes them, and it is the one this program
is aiming at.

%% SOURCE: notes/assets/drawio/perturb-seq-workflow.drawio.svg, exported by
%% `make figures`. Panel numbering and content are unchanged from Fig. 5 of
%% 024-perturb-seq-costing.
\tcfigfit[1.0]{figures/perturb-seq-workflow.pdf}{%
  \textbf{The whole experiment, from library to count matrix, and the one stage
  where the three routes differ.} Read left to right. Two symbols recur in panel
  (b): \textbf{H} is the PCR handle, attached at one end only during barcoding,
  and \textbf{CB} is the cell barcode. The color key is repeated at the foot of
  the figure.
  \textbf{(a)} The stages. Every wet-lab step is common to all three routes
  except the shaded block, where split-pool barcodes cells across three rounds of
  plates, unmodified droplet gives each cell its own bead, and preindexing does
  one plate round first and then a droplet, so several cells can share a droplet
  and still be told apart. Fragmentation is drawn as its own stage for two
  reasons: its position relative to ribosomal RNA depletion decides how much
  that depletion is worth, roughly \$99{,}000 at genome scale, and it determines
  what a count means.
  \textbf{(b)} The same physical cutting, with opposite consequences, which is
  the reason this figure is worth reading rather than skimming. In bulk RNA
  sequencing every piece of a molecule can be sequenced, so a long transcript
  yields proportionally more reads and the length has to be divided back out,
  which is what FPKM and TPM exist to do. In a 3$'$-counting single-cell assay
  only the barcode-proximal piece keeps the outer handle, so exactly one
  fragment survives per molecule whatever its length. Applying TPM or FPKM to a
  count matrix from either route is therefore an error rather than a stylistic
  choice.
  \textbf{(c)} The funnel, which is what a budget is actually made of. Cells are
  lost at recovery, at quality control and at guide assignment, and reads are
  lost separately to the spike-in, to ribosomal sequence and to cells that are
  discarded later. Composed from the protocols rather than traced from a
  published figure, and shared with the costing review, where the numbers drawn
  in it are sourced individually.}{fig:workflow}
%% Source notes/assets/drawio/perturb-seq-workflow.drawio, exported by `make figures`

\notesec{Droplet, which the core runs as a service} A microfluidic chip
co-encapsulates one cell with one barcoded bead, and every molecule released
inside that droplet picks up the bead's barcode. The Carver Center owns the
instrument, the GEM-X 5$'$ kit and the quality control, so our bench work
ends when fixed cells are delivered. Yeast needs one modification, roughly
a microlitre of zymolyase in the reaction to digest the cell
wall~\citep{jarianiNewProtocolSinglecell2020}, and the core has said it accepts
that. A channel recovers about 20{,}000 cells and costs \$2{,}650 for the first
library and \$2{,}240 for each of the next six.

\notesec{Plate, which we would run ourselves} Cells are fixed and permeabilized,
then distributed across a 96-well plate where each well appends a different
barcode; the wells are pooled, re-split and barcoded again, twice more. A cell's
identity is the sequence of wells it visited, so no instrument ever has to
isolate it~\citep{brettnerUltraHighthroughputMassively2024}. One pass through
the protocol handles up to 96 samples and, at a full plate, roughly 480{,}000
cells. The reagents are cheap and the labor is not: about a week of skilled
bench work per run, against a channel the core runs.

\notesec{Preindexed droplet, which is one plate round and then a droplet}
\label{sec:preindexed}
The third route composes the other two, and it is the one worth aiming at.
Permeabilized cells are barcoded once on a 384-well plate, where a well-specific
primer reverse transcribes inside the still-intact cell and writes round one into
the complementary DNA. The cells are then pooled and deliberately overloaded into
a single droplet channel, where a bead writes round
two~\citep{datlingerUltrahighthroughputSinglecellRNA2021}. Neither round
identifies a cell on its own, and that is the entire trick: a droplet may hold
several cells, because the plate already told them apart.

What that removes is the constraint that makes droplet screening expensive. An
ordinary channel runs at low occupancy so that two cells rarely share a droplet,
which leaves most droplets empty and most of the reagent unused. Datlinger et
al.\ loaded 100-fold over the recommendation, filled 95.5\% of droplets at an
average of 9.6 nuclei each, and recovered \textbf{151{,}788} transcriptomes from
one channel, which is 7.6 times the 20{,}000 the core quotes. Since reagent is
sold per channel, a genome-scale screen falls from 137 channels and \$364{,}000
to 18 channels and \$109{,}000. Depleted split-pool is still cheaper at
\$57{,}500, but a factor of 1.9 is small enough to be decided by things the
dollar column does not hold, and all of them favor this route: the core runs the
droplet platform as a service and owns its quality control, the 384-well plate
labels growth conditions for nothing, and an accurate cell count matters far less
when the channel is saturated on purpose.

\textbf{The catch is that it has never been run in any microbe.} Every
demonstration is human or mouse. It also inherits the plate route's sample
preparation rather than the droplet route's: writing a barcode inside a cell
needs reagents to reach the transcriptome in situ, which in yeast means fixation
and wall digestion rather than the in-droplet zymolyase above.
\textbf{Hypothesis (untested):} the two halves compose, since in-cell reverse
transcription on fixed yeast and droplet capture of yeast are each established
separately.

\notesec{A first barcode round on a plate is what makes growth media cheap}
Wherever round one is a plate of wells, its index is read out of every cell
whether or not it has been given a meaning. Growing conditions in a deep-well
plate and barcoding them into the same run therefore costs no barcode space at
all (\cref{fig:pilots}b), and this is the property the two plate-first routes
share: split-pool carries 96 such labels and preindexing carries 384. Only the
unmodified droplet route lacks them, so there each condition buys its own
channel. That is what makes a plate-first route the one to use when the
experiment is about growth media, and it is why the \org{P.\ kudriavzevii} run
should not be a single condition.

\notesec{The guide readout, which is the whole \org{S.\ cerevisiae} question}
\label{sec:guide}
Guide RNAs are transcribed by RNA polymerase III, so they are neither capped nor
polyadenylated, and an assay that captures messenger RNA by its poly(A) tail
sees nothing. That is why our library is read today by amplifying the plasmid
rather than the RNA. The 5$'$ kit may not have the problem: it captures the
guide with a primer against the guide's own scaffold, carried in the
reverse-transcription mix, and 10x state the requirement as an architectural one,
that guides be ``engineered for use with standard Cas9 systems with a
protospacer on the 5$'$ end''\external{10x CG000735}. A conventional polymerase
III cassette, including ours, has exactly that architecture.

\textbf{Hypothesis (untested):} the scaffold primer reverse transcribes a yeast
polymerase III guide, and perturbation identity is therefore readable from an
unmodified library. Two things have to hold and neither is established. The
primer set is proprietary and validated against the standard \gene{SpCas9}
scaffold, and 10x say plainly that ``alternate sgRNA structures \dots\ have not
been tested''\external{10x CG000735}, so our scaffold has to match theirs. And
the poly-dT half of that primer mix contributes nothing to a guide with no tail,
so capture rests entirely on the scaffold primer. If it fails, the fallback is a
polyadenylated guide barcode transcribed by polymerase II, which is a
construction project of several months.

The guide library is sequenced far more shallowly than the transcriptome, 5{,}000
read pairs per cell against 20{,}000, and 10x warn against sequencing it
alone\external{10x CG000735}, so it is pooled with the expression library rather
than submitted separately. It is in the price in \cref{tab:pilots}.

\begin{keybox}
\textbf{One channel, one existing library, and the answer is a number.} The
\org{S.\ cerevisiae} run is worth starting before anything else because it is
the cheapest way to find out whether the whole program needs a new construct.
If $q$ comes back near the 0.71 a published yeast screen reports, single-guide
screening works today and multiplexing is merely expensive. If it comes back
near zero, the 5$'$ scaffold primer does not engage a polymerase III transcript
and a polyadenylated guide barcode has to be built before any screen is
designed.
\end{keybox}
